\section{Compito Tre: l'algoritmo CamminoMin}
\label{sec:compito3}

\subsection{Idea e pseudocodice di riferimento}

Se $D$ appartiene alla chiusura di $O$, un cammino minimo è il cammino libero
fra le due celle e la sua lunghezza è $\dlib(O,D)$. Altrimenti, qualunque
cammino da $O$ a $D$ esce dalla chiusura di $O$ passando per una cella della
sua frontiera. La traccia definisce quindi ricorsivamente il cammino minimo
come il più corto fra i cammini ottenuti componendo il cammino libero da $O$
a una cella di frontiera $F$ con un cammino minimo da $F$ a $D$, al variare di
$F$ nella frontiera. Nella chiamata su $F$ le celle della chiusura di $O$
sono trattate come ostacoli, così il cammino non ripassa da una regione già
coperta.

L'\cref{alg:traccia} riporta lo pseudocodice della traccia (pag.~30); la
numerazione delle righe è quella usata nel resto della relazione.

\begin{algorithm}[H]
\caption{\cmin{} come definito nella traccia~\cite{traccia}}
\label{alg:traccia}
\begin{algorithmic}[1]
\Procedure{CamminoMin}{$O, D, \mathit{dim}, \mathit{osta}$}
  \LineComment{$\mathit{dim}$: dimensioni della griglia; $\mathit{osta}$: celle non attraversabili}
  \State $\mathit{contesto} \gets \Call{CalcolaContesto}{O, \mathit{dim}, \mathit{osta}}$
  \If{$D \in \mathit{contesto}$}
    \State \Return $\dlib(O,D), \seq{(O,0), (D,1)}$
  \EndIf
  \State $\mathit{complemento} \gets \Call{CalcolaComplemento}{O, \mathit{dim}, \mathit{osta}}$
  \If{$D \in \mathit{complemento}$}
    \State \Return $\dlib(O,D), \seq{(O,0), (D,2)}$
  \EndIf
  \State $\mathit{frontiera} \gets \Call{CalcolaFrontiera}{\mathit{contesto}, \mathit{complemento}, \mathit{dim}, \mathit{osta}}$
  \If{$\mathit{frontiera} = \emptyset$} \Comment{vicolo cieco}
    \State \Return $\infty, \seq{}$
  \EndIf
  \State $\mathit{lunghezzaMin} \gets \infty$
  \State $\mathit{seqMin} \gets \seq{}$
  \ForEach{$(F,t) \in \mathit{frontiera}$}
    \State $\mathit{lF} \gets \dlib(O,F)$
    \If{$\mathit{lF} < \mathit{lunghezzaMin}$}
      \LineComment{una condizione più forte è $\mathit{lF} + \dlib(F,D) < \mathit{lunghezzaMin}$}
      \State $\mathit{lFD}, \mathit{seqFD} \gets \Call{CamminoMin}{F, D, \mathit{dim}, \mathit{osta} \cup \mathit{contesto} \cup \mathit{complemento}}$
      \State $\mathit{lTot} \gets \mathit{lF} + \mathit{lFD}$
      \If{$\mathit{lTot} < \mathit{lunghezzaMin}$}
        \State $\mathit{lunghezzaMin} \gets \mathit{lTot}$
        \State $\mathit{seqMin} \gets \Call{Compatta}{\seq{(O,0),(F,t)}, \mathit{seqFD}}$
      \EndIf
    \EndIf
  \EndFor
  \State \Return $\mathit{lunghezzaMin}, \mathit{seqMin}$
\EndProcedure
\end{algorithmic}
\end{algorithm}

\subsection{Implementazione}
\label{sec:implementazione-cmin}

La classe \code{CamminoMin} implementa l'\cref{alg:impl}. Il metodo
pubblico \code{calcola(griglia, O, D)} restituisce un
\code{Cammino}, formato dalla lunghezza (\code{DistanzaLibera}) e dalla lista
dei landmark. Due sovraccarichi ricevono anche un oggetto
\code{StatisticheEsecuzione}, in cui accumulare i contatori del riassunto, e la
variante dell'algoritmo da eseguire (\cref{sec:varianti}); senza indicazioni
si usa la variante \code{base}.

\begin{algorithm}[tbp]
\caption{\cmin{} come implementato in \code{CamminoMin}}
\label{alg:impl}
\begin{algorithmic}[1]
\Require griglia $g$, celle attraversabili $O$ e $D$, limite $L$ ($\infty$ nella chiamata principale), tabella $\mathit{cache}$ (vuota alla prima chiamata)
\Ensure lunghezza del cammino minimo da $O$ a $D$ e sequenza dei landmark, oppure $(\infty, \seq{})$ se non esiste un cammino più corto di $L$
\Function{CamminoMin}{$g, O, D, L$}
  \State $k \gets \Call{Chiave}{\text{ostacoli di } g, O, D}$
  \If{$\mathit{cache}[k]$ è esatto, o esclude cammini più corti di un $L' \ge L$}
    \State \Return $\mathit{cache}[k]$ \Comment{$(\infty, \seq{})$ nel secondo caso}
  \EndIf
  \State $s, \mathit{chiusura}, \mathit{frontiera} \gets \Call{CalcolaChiusura}{g, O}$ \Comment{\cref{alg:chiusura}}
  \If{$s[D]$ contiene \code{CHIUSURA}} \Comment{righe 4--8, anche $O = D$}
    \State $r \gets \bigl(\dlib(O,D), \seq{(O,0), (D, \Call{Tipo}{s[D]})}\bigr)$
    \State $\mathit{cache}[k] \gets r$ (esatto); \quad \Return $r$
  \EndIf
  \State $\mathit{lunghezzaMin} \gets L$; \quad $\mathit{seqMin} \gets \seq{}$
  \State ordina $\mathit{frontiera}$ per $\dlib(O,F) + \dlib(F,D)$ crescente
  \ForEach{$F \in \mathit{frontiera}$}
    \State $\mathit{lF} \gets \dlib(O,F)$
    \If{$\mathit{lF} \ge \mathit{lunghezzaMin}$} \Continue \Comment{riga 16 falsa}
    \EndIf
    \If{$\mathit{lF} + \dlib(F,D) \ge \mathit{lunghezzaMin}$} \Continue \Comment{riga 17 falsa}
    \EndIf
    \State $g' \gets g$ con le celle di $\mathit{chiusura} \setminus \{F\}$ come ostacoli
    \State $L' \gets \mathit{lunghezzaMin} - \mathit{lF}$ se il limite è propagato, altrimenti $\infty$
    \State $\mathit{lFD}, \mathit{seqFD} \gets \Call{CamminoMin}{g', F, D, L'}$
    \If{$\mathit{lF} + \mathit{lFD} < \mathit{lunghezzaMin}$}
      \State $\mathit{lunghezzaMin} \gets \mathit{lF} + \mathit{lFD}$
      \State $\mathit{seqMin} \gets \seq{(O,0), (F, \Call{Tipo}{s[F]})} \cdot \mathit{seqFD}[2..]$ \Comment{\textsc{Compatta}}
    \EndIf
  \EndFor
  \If{$\mathit{seqMin} = \seq{}$} \Comment{nessun cammino più corto di $L$}
    \State $\mathit{cache}[k] \gets$ ``nessun cammino più corto di $L$''; \quad \Return $(\infty, \seq{})$
  \EndIf
  \State $\mathit{cache}[k] \gets (\mathit{lunghezzaMin}, \mathit{seqMin})$ (esatto)
  \State \Return $\mathit{lunghezzaMin}, \mathit{seqMin}$
\EndFunction
\end{algorithmic}
\end{algorithm}

Senza limite propagato $L$ vale sempre $\infty$ e l'algoritmo coincide con la
versione della traccia arricchita da ordinamento, riga~17 e cache.

La \cref{tab:corrispondenza} collega le righe della traccia
all'implementazione. La traccia ammette una codifica diversa dallo
pseudocodice purché produca gli stessi effetti (pag.~34--35).

\begin{table}[H]
\centering
\small
\begin{tabularx}{\textwidth}{@{}l X@{}}
\toprule
\textbf{Righe della traccia} & \textbf{Implementazione} \\
\midrule
3, 6, 9  & Un'unica chiamata a \code{CalcoloChiusura} produce contesto,
           complemento, frontiera e tipo di ogni cella (\cref{sec:chiusura}). \\
4--8     & Il caso base controlla la maschera \code{CHIUSURA} nello stato di
           $D$; il tipo di $D$ è $1$ se ha il bit \code{CONTESTO}, $2$
           altrimenti. Se $O = D$, $O$ è nel proprio contesto e il risultato è
           $\bigl(0, \seq{(O,0),(O,1)}\bigr)$, come richiesto a pag.~32. \\
10--11   & Con frontiera vuota il ciclo non esegue iterazioni e la funzione
           restituisce $(\infty, \seq{})$, cioè lo stesso risultato. \\
12--13   & $\mathit{lunghezzaMin}$ parte da $\infty$, oppure dal limite
           ricevuto se il limite è propagato (\cref{sec:ottimizzazioni}). \\
14       & La frontiera è visitata in ordine di limite inferiore crescente.
           $O$ non è mai di frontiera: i suoi vicini attraversabili sono le
           prime celle delle linee primarie e appartengono al contesto. \\
16--17   & La condizione della riga 16 è controllata per prima, poi quella
           più forte suggerita alla riga 17. I due esiti negativi sono
           contati separatamente (\cref{sec:contatori}). \\
18       & Gli ostacoli della chiamata ricorsiva sono quelli di $g$ più la
           chiusura di $O$ tranne $F$ (\cref{sec:ostacoli-annidati}). \\
20--22   & Il confronto è stretto: a parità di lunghezza resta il primo
           cammino trovato. \textsc{Compatta} sostituisce il primo landmark
           $(F,0)$ di $\mathit{seqFD}$ con $(F,t)$. \\
\bottomrule
\end{tabularx}
\caption{Corrispondenza fra lo pseudocodice della traccia e l'implementazione.}
\label{tab:corrispondenza}
\end{table}

\subsection{Ostacoli nelle chiamate annidate}
\label{sec:ostacoli-annidati}

Secondo la traccia (pag.~28 e 32), nella chiamata su $F$ tutte le celle della
chiusura di $O$ diventano ostacoli, compresa $F$, che però resta l'origine
della nuova chiusura. L'implementazione ottiene lo stesso effetto escludendo
$F$ dall'insieme di ostacoli aggiunti: le celle raggiungibili da $F$ sono le
stesse, e \code{CalcoloChiusura} non deve gestire un'origine marcata come
ostacolo. La nuova griglia $g'$ è una vista immutabile
(\cref{sec:immutabilita}): ogni ramo ricorsivo accumula i propri ostacoli
senza modificare la griglia dei rami fratelli.

\subsection{Sequenza dei landmark}

La sequenza dei landmark $\seq{(O,0), (F_1,t_1), \dots, (D,t_D)}$ elenca
l'origine e le celle di frontiera attraversate dal cammino; il tipo di un
landmark è quello del cammino libero che lo raggiunge dal landmark
precedente, quindi la sequenza descrive il cammino in modo univoco.
Nell'implementazione un landmark è un record con le coordinate e il tipo
$t \in \{0, 1, 2\}$; il tipo $0$ identifica l'origine della chiamata. La
lista dei landmark di un \code{Cammino} è immutabile. Se $D$ è
irraggiungibile il risultato è la lunghezza infinita con la sequenza vuota.

Nell'esempio della \cref{fig:esempio}(b) la chiamata da $O=(0,0)$ restituisce
$\seq{(O,0), ((8,7),1), (D,1)}$. La chiamata nel verso opposto, da
$(10,10)$ a $(0,0)$, restituisce
$\seq{((10,10),0), ((0,1),2), ((0,0),1)}$: un cammino diverso, ma della stessa
lunghezza $9\sqrt{2}+2$. In entrambi i casi la somma delle distanze libere
fra landmark consecutivi è uguale alla lunghezza del cammino; i test
verificano questa proprietà su tutti i risultati finiti.

\subsection{Ottimizzazioni}
\label{sec:ottimizzazioni}

Le ottimizzazioni riducono il numero di chiamate ricorsive senza cambiare la
lunghezza restituita; il loro effetto è misurato nella \cref{sec:ablation}.

\paragraph{Ordinamento della frontiera.} Le celle di frontiera sono visitate
in ordine crescente di $\dlib(O,F) + \dlib(F,D)$, un limite inferiore della
lunghezza di qualunque cammino da $O$ a $D$ che passi per $F$. I candidati
più promettenti vengono esplorati per primi, $\mathit{lunghezzaMin}$ scende
presto e le condizioni di potatura scartano più rami. L'ordinamento costa
$O(f \log f)$, con $f$ numero di celle di frontiera; le due distanze libere di
ogni candidato sono calcolate una volta sola e riusate dalla potatura.

\paragraph{Condizione della riga 17.} Dopo la riga 16 viene controllata la
condizione più forte $\mathit{lF} + \dlib(F,D) < \mathit{lunghezzaMin}$. La
potatura è corretta perché $\dlib(F,D)$ non supera la lunghezza di nessun
cammino da $F$ a $D$: se la condizione è falsa, passare per $F$ non può dare
un cammino strettamente più corto di quello già trovato. Il confronto
sperimentale fra le due condizioni, facoltativo per i gruppi di tre persone
(pag.~67), non fa parte di questo progetto.

\paragraph{Memoizzazione.} Il risultato di una chiamata dipende dagli
ostacoli della griglia, da $O$, da $D$ e, con il limite propagato, dal limite
ricevuto. Il risultato viene salvato in una \code{HashMap} la cui chiave è
formata dal \code{BitSet} degli ostacoli della vista (\cref{sec:immutabilita})
e dalle coordinate di $O$ e $D$: la chiave occupa $n/8$ byte e
\code{equals} e \code{hashCode} costano $O(n/64)$. La chiave è esatta, perché
un hash intero degli ostacoli potrebbe collidere e restituire il risultato di
un altro sottoproblema. Una prima versione usava una stringa con le
coordinate di tutti gli ostacoli: chiavi equivalenti, ma costruite con una
scansione della griglia e lunghe decine di caratteri per ostacolo. Le
direzioni $O \to D$ e $D \to O$ usano tabelle separate, così la verifica di
simmetria confronta due calcoli indipendenti.

\paragraph{Limite propagato.} Nella versione della traccia ogni chiamata
annidata cerca il cammino minimo da $F$ a $D$ partendo da
$\mathit{lunghezzaMin} = \infty$, anche quando il chiamante ha già un cammino
di lunghezza $\mathit{lunghezzaMin}$ e gli serve soltanto sapere se ne esiste
uno più corto di $\mathit{lunghezzaMin} - \dlib(O,F)$. La variante con limite
propagato passa questo valore $L'$ alla chiamata annidata, che inizializza
con esso il proprio $\mathit{lunghezzaMin}$: le righe 16 e 17 potano così
anche all'interno della ricorsione. Una chiamata con limite $L$ restituisce
il cammino minimo del sottoproblema se è più corto di $L$, altrimenti
$(\infty, \seq{})$. Il risultato resta corretto: un cammino trovato sotto il
limite viene confrontato con tutti i candidati che potrebbero batterlo, ed è
quindi il minimo; un ramo che non trova nulla sotto il limite non può
migliorare il cammino del chiamante. La cache distingue i due casi: un
risultato trovato è esatto e vale per qualunque limite; un risultato vuoto
significa soltanto ``nessun cammino più corto di $L$'' e viene riusato solo
per chiamate con limite non superiore a $L$.

\paragraph{Viste immutabili.} La griglia $g'$ della chiamata ricorsiva
richiede la copia di un \code{BitSet} di $n$ bit e l'inserimento delle celle
della chiusura, invece della copia di una matrice di $n$ interi.

\subsection{Varianti}
\label{sec:varianti}

L'enumerazione \code{VarianteCamminoMin} attiva o disattiva le
ottimizzazioni, per confrontarle nella sperimentazione
(\cref{tab:varianti}). La riga~17 è attiva in tutte le varianti. La variante
si sceglie da riga di comando con \code{-{}-variante}.

\begin{table}[H]
\centering
\small
\begin{tabular}{@{}l c c c@{}}
\toprule
\textbf{Variante} & \textbf{Cache} & \textbf{Ordinamento} & \textbf{Limite propagato} \\
\midrule
\code{base}              & sì & sì & no \\
\code{senza-cache}       & no & sì & no \\
\code{senza-ordinamento} & sì & no, ordine di scansione della griglia & no \\
\code{limite-propagato}  & sì & sì & sì \\
\bottomrule
\end{tabular}
\caption{Varianti di \cmin{}. Tutte restituiscono la stessa lunghezza minima.}
\label{tab:varianti}
\end{table}

\subsection{Contatori e riassunto}
\label{sec:contatori}

La traccia chiede, per ogni istanza e su richiesta dell'utente, un riassunto
con dimensioni e tipo della griglia, numero totale di celle di frontiera,
numero di volte in cui la condizione della riga 16/17 è risultata falsa e
prestazioni (pag.~71). \code{StatisticheEsecuzione} raccoglie i contatori
della \cref{tab:contatori} durante l'esecuzione. Un'istanza comprende sempre
entrambe le direzioni $O \to D$ e $D \to O$; i contatori sommano le due
direzioni e tutte le chiamate ricorsive.

\begin{table}[H]
\centering
\small
\begin{tabularx}{\textwidth}{@{}l X@{}}
\toprule
\textbf{Campo} & \textbf{Significato} \\
\midrule
\code{front}         & celle di frontiera di tutte le chiusure calcolate
                       (le chiamate servite dalla cache non calcolano chiusure) \\
\code{r16False}      & volte in cui la condizione della riga 16 è risultata falsa \\
\code{r17Prune}      & volte in cui la riga 16 era vera ma la condizione più forte
                       della riga 17 era falsa \\
\code{pruneTot}      & somma dei due precedenti: volte in cui la condizione
                       effettivamente usata è risultata falsa e il ramo è stato potato \\
\code{variante}      & variante di \cmin{} eseguita (\cref{tab:varianti}) \\
\code{chiamate}      & invocazioni di \cmin{}, comprese quelle servite dalla cache \\
\code{cacheHit}      & invocazioni servite dalla cache, che non calcolano una chiusura \\
\code{prof}          & profondità massima della ricorsione \\
\code{chiusurePicco} & massimo numero di matrici di chiusura presenti
                       contemporaneamente in memoria \\
\code{memEst}        & stima della memoria di tali matrici:
                       $\code{chiusurePicco} \cdot 4n$ byte \\
\code{tempoNs}       & tempo minimo, mediano e massimo delle ripetizioni (\cref{sec:protocollo}) \\
\code{seqOD}, \code{seqDO} & sequenze dei landmark nei due versi \\
\bottomrule
\end{tabularx}
\caption{Contatori del riassunto. L'elenco completo delle colonne del CSV è
nell'\cref{app:formato}.}
\label{tab:contatori}
\end{table}

Poiché $\dlib(F,D) \ge 0$, quando la condizione della riga 16 è falsa lo è
anche quella della riga 17. \code{pruneTot} conta quindi le volte in cui la
condizione della riga 17 è risultata falsa, e \code{r16False} quante di
queste sarebbero state potate anche dalla sola riga 16. La profondità e le
chiusure presenti in memoria possono differire: una chiamata servita dalla
cache aumenta la profondità ma non alloca una nuova matrice.

Il riassunto viene stampato su una riga per istanza e variante, con le
lunghezze arrotondate a quattro decimali; nel piano sperimentale i riassunti
sono salvati anche in \code{risultati/riassunti.txt}. Per un'istanza singola
l'opzione \code{-{}-summary false} omette la stampa, mentre il CSV contiene
sempre tutti i campi, con sei decimali. Il riassunto dell'istanza della \cref{fig:esempio} nel piano
sperimentale è:

\begin{lstlisting}
[G12x12_d10_s50_p1] variante=base | tipo=Casuale | 12x12 | ost=15 | O=(0,0)->D=(10,10) | OD=14.7279 | DO=14.7279 | corr=true | front=118 | r16False=1 | r17Prune=96 | pruneTot=97 | chiamate=18 | cacheHit=2 | prof=4 | chiusurePicco=4 | memEst=2304B | tempoNs[min/med/max]=180900/184000/218900 | warmup=1 | rip=5 | seqOD=[(x=0,y=0,tipo=0), (x=8,y=7,tipo=1), (x=10,y=10,tipo=1)] | seqDO=[(x=10,y=10,tipo=0), (x=0,y=1,tipo=2), (x=0,y=0,tipo=1)]
\end{lstlisting}
